\documentclass[review]{elsarticle}
%% review  = one-column, double-spaced, line numbers (typical submission format)
%% replace with \documentclass[final,5p,times]{elsarticle} for the camera-ready two-column layout later.

\usepackage{amsmath,amssymb}
\usepackage{graphicx}
\usepackage{booktabs}
\usepackage{hyperref}

\journal{Computers \& Structures}

\bibliographystyle{elsarticle-num}

\begin{document}

\begin{frontmatter}

\title{Physics-Informed Neural Operator Learning for Hyperelastic Solids: From 2D Benchmarks to a 3D Rocking Rubber-Mount Bushing}

%% TODO: confirm author list and order with Omar/Timon before submission.
\author[rub]{Omar Amro}
\author[rub]{Timon Rabczuk\corref{cor1}}
\cortext[cor1]{Corresponding author.}
\address[rub]{Chair of Computational Mechanics, Bauhaus-Universit\"at Weimar}

\begin{abstract}
%% DRAFT -- placeholder numbers below must be re-checked against
%% PROJECT_STATUS.md immediately before submission; written for review now.
We train a physics-informed neural operator (Transolver) to predict the
displacement field of hyperelastic solids under randomly sampled material
and loading fields, without any labeled displacement data: the network is
trained by directly minimizing the discrete total potential energy of the
body (a Deep Energy Method), following the PFEM pretraining methodology.
Two 2D benchmark geometries (a unit square under top-edge traction, and a
quarter ring under internal pressure), each combined with three isotropic
hyperelastic strain-energy densities (Neo-Hookean, Mooney--Rivlin,
Arruda--Boyce), are trained to completion, evaluated against an independent
finite-element reference, and characterized across mesh-resolution
invariance, out-of-distribution robustness, and break-even cost against both
a native and a GPU-native finite-element solver. The study is then extended
to a third, more demanding 3D benchmark: a rocking elastomeric rubber-mount
bushing with a stress-concentrating groove. The operator reproduces
displacement, strain energy, and reaction moment accurately (combined
displacement error 1.9\%, energy error 0.3\%, reaction-moment error 2.7\%),
and, judged against the cheapest finite-element mesh that matches its own
accuracy, remains 589--800$\times$ faster per sample while repaying its
one-off training cost within fewer than 1{,}900 solves. The one quantity
that does not reach the same standard is the regional Cauchy-stress field
near the groove: an input-normalization fix first resolves a training
instability, after which the regional-stress error is traced, in part, to
an under-resolved background numerical integration rather than the
network's own capacity, and a local integration refinement -- applied to
only the few mesh elements touching the region of interest, independently
of the operator's own discretization -- roughly halves the error (from
62\% to 30\%). The remaining gap is reported as an open limitation, with
several concrete directions (stress/equilibrium-aware physics losses,
mixed displacement--stress formulations, local/hierarchical representations)
identified for future work.
\end{abstract}

\begin{keyword}
physics-informed neural operators \sep deep energy method \sep Transolver \sep hyperelasticity \sep finite element method \sep operator learning
\end{keyword}

\end{frontmatter}

\section{Introduction}
\label{sec:intro}

Classical finite-element (FEM) solvers for nonlinear elasticity re-solve a
new nonlinear system from scratch for every new combination of geometry,
material field, and loading, each solve requiring a full Newton--Raphson
iteration to convergence. Operator learning instead trains a single neural
network to approximate the map from problem data (material field, loading
field, geometry) directly to the solution field, so that new instances are
evaluated with a single forward pass~\cite{lu2021deeponet,li2021fno,kovachki2023neuraloperator}.
Physics-informed variants of this idea remove the need for expensive
labeled training data entirely, by training the network to satisfy the
governing equations or minimize a governing energy functional directly,
rather than to regress against precomputed solutions~\cite{raissi2019pinn,e2018deepritz}.

This work follows the PFEM (\emph{Pretrain Finite Element Method})
approach~\cite{wang2026pfem}: a Transolver-architecture
operator~\cite{wu2024transolver}, built on the attention
mechanism~\cite{vaswani2017attention}, is trained without any labeled
displacement data, by directly minimizing the discretized total potential
energy of the elastic body -- a Deep Energy Method (DEM)
objective~\cite{samaniego2020energy,nguyen2020deep}, closely related to
the Deep Ritz method~\cite{e2018deepritz} and to the variational
physics-informed operator framework~\cite{eshaghi2025vino}, rather than a
supervised regression loss against finite-element solutions. The same
underlying idea -- minimizing an energy functional with a neural ansatz,
instead of imposing the strong-form PDE residual as in standard PINNs --
has previously been developed for finite-deformation
hyperelasticity~\cite{nguyen2020deep,abueidda2022dlem}, strain-gradient
elasticity~\cite{nguyen2021parametric}, contact
problems~\cite{bai2025energy}, and fracture~\cite{goswami2020transfer,wang2026xdem,zhang2026meshfree},
predominantly at the level of a single geometry/material instance rather
than as a trained operator generalizing across a distribution of material
and loading fields.

The present study covers two 2D benchmark geometries -- a unit square
under top-edge traction (B1) and a quarter ring under internal pressure
(B2) -- combined with three hyperelastic strain-energy densities
(Neo-Hookean~\cite{rivlin1948large}, Mooney--Rivlin~\cite{mooney1940theory},
Arruda--Boyce~\cite{arruda1993threedim}), for six total (geometry,
material) cases, and a third, substantially more demanding 3D benchmark: a
rocking elastomeric rubber-mount bushing with a stress-concentrating
groove in its rigid inner core (B3). An independent Newton--Raphson
Total-Lagrangian finite-element solver, built on top of
torch-fem~\cite{meyer2024torchfem} and
PyTorch~\cite{paszke2019pytorch,baydin2018automatic}, generates
ground-truth samples used only for validation, never for training.

%% TODO: contribution list / paper outline paragraph.

\section{Governing Equations and Material Models}
\label{sec:governing}
%% TODO: port from Report Section 2 (Total Lagrangian hyperelasticity;
%% Neo-Hookean / Mooney-Rivlin / Arruda-Boyce strain-energy densities),
%% condensed to paper length.

\section{Benchmark Geometries}
\label{sec:geometries}
%% TODO: port from Report Section 3 (B1, B2) + Section 11.1 (B3 geometry
%% and design history), condensed.

\section{Spatial Discretization and Dataset Generation}
\label{sec:discretization}
%% TODO: port from Report Section 4.

\section{Neural Operator Architecture and Training Objective}
\label{sec:architecture}
%% TODO: port from Report Section 5 (Transolver; DEM training objective).

\section{Results: B1/B2 (2D Benchmarks)}
\label{sec:results-2d}
%% TODO: port from Report Sections 6-10 (training protocol, main results,
%% GPU-native FEM comparison, OOD, resolution invariance, physics-informed
%% vs. data-driven, verification against manufactured solutions).

\section{Results: B3 (3D Rocking Rubber-Mount Bushing)}
\label{sec:results-3d}
%% TODO: port from Report Section 11.6-11.10 (dataset/training,
%% displacement/energy/reaction accuracy, the regional Cauchy-stress
%% investigation and local integration refinement, the accuracy-matched
%% break-even analysis), condensed to paper length and reframed from
%% "report" register to "paper" register (no process narrative, results
%% and methodology only).

\section{Discussion}
\label{sec:discussion}
%% TODO: port from Report Section 9, plus the B3 regional-stress
%% limitation and its future-work directions (physics/equilibrium-aware
%% losses, mixed formulations, PINO-type residual enforcement,
%% hierarchical representations, data-driven stress supervision, hybrid
%% approaches; separating physics-integration resolution from operator
%% resolution more generally) -- citing
%% \cite{abueidda2023enhanced,rao2021elastodynamics,abueidda2021meshless,he2023gnn,taghikhani2025neuralnewton,dean2026piginot}
%% as related attempts, and \cite{prudhomme1999goal,becker2001optimal}
%% for the goal-oriented-error-estimation framing of "where does a
%% trained operator's error actually come from."

\section{Conclusion}
\label{sec:conclusion}
%% TODO: port from Report Section 10, extended with the B3 conclusion.

\section*{Acknowledgments}
%% TODO.

\bibliography{references}

\end{document}
